\begin{table}[!htbp]
\caption{Shock-Level Diagnostics for the Composition Instrument}
\label{tab:bhj_shock_level}
{\centering
\adjustbox{max width=\linewidth}{%
\begin{tabular}{lr}
\toprule
Shock-level observations $(k,t)$ & 9,600 \\\\
Effective shocks, each $(k,t)$ & 1108.1 \\\\
Effective shocks, each category once & 40.0 \\\\
County-level 2SLS coefficient & 4.0909 \\\\
Shock-level IV coefficient & 4.0909 \\\\
Equivalence gap & 1.4e-15 \\\\
Coefficient, no share-sum control & 4.091 \\\\
\hspace{1.5em}$t$ & 6.77 \\\\
Coefficient, share sum controlled & 3.367 \\\\
\hspace{1.5em}$t$ & 5.46 \\\\
Share sum's share of instrument variance & 0.052 \\\\
Balance: income per head & -0.0071 (0.0183) \\\\
Balance: employment rate & -0.0010 (0.0154) \\\\
Balance: log population & 0.0034 (0.0103) \\\\
Balance: pre-period growth & 0.0114 (0.0209) \\\\
Balance: joint $F$ & 0.59 \\\\
Observations, county level & 97,856 \\\\
\bottomrule
\end{tabular}}
\par}
\medskip
{\scriptsize \noindent\textit{Notes:} The composition instrument at a four-year horizon on
97,856 county-years, with division-by-year fixed effects, division-equal weights and the
treatment winsorized within year. Following Borusyak, Hull and Jaravel, the shift-share IV is
rewritten as an IV at the level of the 300 Federal Supply Classification categories,
using exposure-weighted aggregates of the outcome and the treatment and weighting by shock
importance; the two coefficients agree to the reported gap, so the design is a shock-level
experiment in their sense. Effective shock counts are inverse Herfindahl indices of the
importance weights, reported both treating each category-year as a shock and treating each
category as one shock, the latter being conservative because a category's national series is
serially correlated. Their framework requires exposure shares summing to a constant; here they
sum to total contract dollars per head, so the estimate is reported with and without that sum
partialled out. Balance regressions take the national category rate on exposure-weighted
predetermined county characteristics, standardized, weighted by shock importance and clustered
by category; the pre-period growth rate is the same midpoint employment growth measured over
the four years ending where the estimation window begins. Balance failing to reject is evidence
consistent with quasi-random assignment and is not proof of it.\par}
\end{table}
